Este artículo presenta un estudio de investigación aplicada sobre la utilización de Arquitecturas Orientadas a Eventos (EDA) y microservicios como remplazo a arquitecturas obsoletas en el desarrollo de software moderno. La investigación se centra específicamente en ecosistemas de Internet de las Cosas (IoT), debido a que en este tipo de ecosistemas el rendimiento y el procesamiento de datos en tiempo real a grandes velocidades es vital, lo que provoca que el uso de arquitecturas antiguas, genere un sistema obsoleto debido a su bajo rendimiento en el procesamiento continuo de datos y no solo se quedan atrás en tema de rendimiento y procesamiento de datos, sino que los sensores carecen de contexto y de cualquier forma de comunicación externa, lo que impide la comunicación entre sensores y el procesamiento de los datos se hace directamente en la nube, aumentando así los gastos en soluciones en la nube para la mejora del procesamiento de datos y el rendimiento. En este articulo analizamos varias soluciones a estas amenazas, incluyendo arquitecturas orientadas a eventos basadas en microservicios, arquitecturas orientadas a eventos implementando Computación en el Borde (Edge Computing) agregándole a cada uno de los sensores su propio contexto y la posibilidad de comunicarse con otros sensores para la creación incluso de eventos complejos, además de mejorar el rendimiento en el procesamiento de datos, y aún mejor, solucionar uno de los mayores problemas de los sensores IoT que es la normalización de datos heterogéneos, Asimismo, se evalúa el impacto de utilizar mensajería centralizada mediante un bus de eventos (Apache Kafka) para garantizar el desacoplamiento y la comunicación asíncrona entre componentes. Los resultados que se obtuvieron demuestran una mejora significativa en el rendimiento y eficiencia en los sistemas IoT cuando adoptaron arquitecturas orientadas a eventos y microservicios, en comparación con las arquitecturas tradicionales. También se observó una reducción notable en los costos relacionados con el procesamiento de datos en la nube.